% titel: Blatt 9, A6: Potenzreihe für ein Integral (Gauß-Integral)
% Teil II: der durchgerechnete Weg.

\subsection*{Gegeben und gesucht}
(a) Eine Potenzreihe $\sum_{n=0}^{\infty}a_n x^n$, die für alle $x$ konvergiert
und gleich $\int_0^x \mathrm e^{-t^2/2}\,\dd t$ ist, analog Beispiel 5.4.5,3.
(b) Python-Code, der $\int_0^1 \mathrm e^{-t^2/2}\,\dd t$ mittels dieser Reihe
approximiert.

\subsection*{a) Potenzreihe für $\int_0^x \mathrm e^{-t^2/2}\,\dd t$}
Die Setupreihe der Exponentialfunktion gibt für alle $t$:
\begin{equation}
  \mathrm e^{-t^2/2}
  = \sum_{k=0}^{\infty}\frac{(-t^2/2)^k}{k!}
  = \sum_{k=0}^{\infty} (-1)^k \frac{t^{2k}}{2^k\,k!}.
\end{equation}
Gliedweises Integrieren von $0$ bis $x$ (Satz 5.4.2; Konvergenzradius bleibt
$\infty$):
\begin{align}
  \int_0^x \mathrm e^{-t^2/2}\,\dd t
  &= \sum_{k=0}^{\infty} (-1)^k \frac{1}{2^k\,k!}\int_0^x t^{2k}\,\dd t \\
  &= \sum_{k=0}^{\infty} (-1)^k \frac{x^{2k+1}}{(2k+1)\,2^k\,k!}.
\end{align}
Dies ist eine Potenzreihe $\sum_n a_n x^n$ mit
$a_{2k+1}=\dfrac{(-1)^k}{(2k+1)\,2^k\,k!}$ und $a_{2k}=0$.

\begin{ergebnis}
$\displaystyle
\int_0^x \mathrm e^{-t^2/2}\,\dd t
=\sum_{k=0}^{\infty}(-1)^k\frac{x^{2k+1}}{(2k+1)\,2^k\,k!},\qquad
a_{2k+1}=\frac{(-1)^k}{(2k+1)\,2^k\,k!},\ a_{2k}=0$.
\end{ergebnis}

\subsection*{b) Python-Approximation von $\int_0^1 \mathrm e^{-t^2/2}\,\dd t$}
Für $x=1$ in die Reihe einsetzen und gliedweise aufsummieren:
\begin{verbatim}
import math
N = 50
s = 0.0
for k in range(N):
    s += (-1)**k / ((2*k+1) * 2**k * math.factorial(k))
print(s)
\end{verbatim}
Analog lösbar über die rekursive Glieddefinition
`term *= -x**2 * (2*k+1) / (2*(k+1))` in einer Schleife. Für $N=50$ erhält man
$s\approx0.85562439$.

\begin{ergebnis}
$\displaystyle\int_0^1 \mathrm e^{-t^2/2}\,\dd t \approx 0.85562439 .$
\end{ergebnis}

\begin{kontrolle}
\textbf{Referenz:} Das Integral ist $\sqrt{\tfrac\pi2}\,\mathrm{erf}(\tfrac1{\sqrt2})
\approx0.85562439$ — stimmt mit der Reihensumme überein.
\textbf{Grenzfall:} Integrand $\mathrm e^{-t^2/2}$ ist positiv, nie größer als
$1$; das Integral über $[0,1]$ liegt zwischen $0$ und $1$, das Ergebnis
$0.856$ ist plausibel.
\end{kontrolle}

\begin{merksatz}
Eine Funktion ohne elementare Stammfunktion (hier $\mathrm e^{-t^2/2}$) wird
über die Potenzreihe integrierbar: Reihe hinschreiben, gliedweise integrieren,
Danach bleibt der Konvergenzradius erhalten. Für numerische Werte reicht das
Partial aufsummieren der alternierenden Reihe.
\end{merksatz}
